\chapter{Unit 4: Fundamentals of Semiconductor Devices}

\section{Section A: Multiple Choice Question Practice (60 Curated Problems)}

\noindent\textbf{Q1. Which of the following is the most commonly used semiconductor material in electronic devices?} \hfill [Topic: introduction to semiconductor devices | Easy]
\begin{enumerate}[label=(\alph*),topsep=2pt,itemsep=1pt,partopsep=0pt,parsep=0pt]
    \item {Aluminium}
    \item {Silicon}
    \item {Copper}
    \item {Mica}
\end{enumerate}
\vspace{2mm}

\noindent\textbf{Q2. A pure semiconductor without any impurity added is known as a:} \hfill [Topic: introduction to semiconductor devices | Easy]
\begin{enumerate}[label=(\alph*),topsep=2pt,itemsep=1pt,partopsep=0pt,parsep=0pt]
    \item {P-type semiconductor}
    \item {Extrinsic semiconductor}
    \item {N-type semiconductor}
    \item {Intrinsic semiconductor}
\end{enumerate}
\vspace{2mm}

\noindent\textbf{Q3. In an N-type semiconductor, the majority charge carriers are:} \hfill [Topic: introduction to semiconductor devices | Easy]
\begin{enumerate}[label=(\alph*),topsep=2pt,itemsep=1pt,partopsep=0pt,parsep=0pt]
    \item {Protons}
    \item {Neutrons}
    \item {Electrons}
    \item {Holes}
\end{enumerate}
\vspace{2mm}

\noindent\textbf{Q4. When a PN junction diode is forward biased, the depletion region:} \hfill [Topic: Working and operation of PN Junction diode | Easy]
\begin{enumerate}[label=(\alph*),topsep=2pt,itemsep=1pt,partopsep=0pt,parsep=0pt]
    \item {Decreases in width}
    \item {Increases in width}
    \item {Becomes infinitely wide}
    \item {Stays exactly the same}
\end{enumerate}
\vspace{2mm}

\noindent\textbf{Q5. A PN junction diode primarily allows current to flow in:} \hfill [Topic: Working and operation of PN Junction diode | Easy]
\begin{enumerate}[label=(\alph*),topsep=2pt,itemsep=1pt,partopsep=0pt,parsep=0pt]
    \item {No direction at all}
    \item {Both directions equally}
    \item {Only one direction}
    \item {Only alternating directions}
\end{enumerate}
\vspace{2mm}

\noindent\textbf{Q6. The approximate forward voltage drop (knee voltage) of a silicon PN junction diode is about:} \hfill [Topic: Characteristics of PN junction diode | Easy]
\begin{enumerate}[label=(\alph*),topsep=2pt,itemsep=1pt,partopsep=0pt,parsep=0pt]
    \item {1.5 V}
    \item {5.0 V}
    \item {0.3 V}
    \item {0.7 V}
\end{enumerate}
\vspace{2mm}

\noindent\textbf{Q7. In the reverse-biased region of an ideal PN junction diode, the current is:} \hfill [Topic: Characteristics of PN junction diode | Easy]
\begin{enumerate}[label=(\alph*),topsep=2pt,itemsep=1pt,partopsep=0pt,parsep=0pt]
    \item {Equal to forward current}
    \item {Nearly zero}
    \item {Very large}
    \item {Infinite}
\end{enumerate}
\vspace{2mm}

\noindent\textbf{Q8. A Zener diode is normally operated in which region?} \hfill [Topic: working and operation of Zener diode | Easy]
\begin{enumerate}[label=(\alph*),topsep=2pt,itemsep=1pt,partopsep=0pt,parsep=0pt]
    \item {Reverse breakdown region}
    \item {Cut-off region only}
    \item {Saturation region}
    \item {Forward conduction region}
\end{enumerate}
\vspace{2mm}

\noindent\textbf{Q9. The main application of a Zener diode is as a:} \hfill [Topic: working and operation of Zener diode | Easy]
\begin{enumerate}[label=(\alph*),topsep=2pt,itemsep=1pt,partopsep=0pt,parsep=0pt]
    \item {Current amplifier}
    \item {Signal oscillator}
    \item {Voltage regulator}
    \item {Light emitter}
\end{enumerate}
\vspace{2mm}

\noindent\textbf{Q10. In the Zener breakdown region, the voltage across the Zener diode remains:} \hfill [Topic: characteristics of Zener diode | Easy]
\begin{enumerate}[label=(\alph*),topsep=2pt,itemsep=1pt,partopsep=0pt,parsep=0pt]
    \item {Rapidly increasing}
    \item {Steadily decreasing}
    \item {Equal to zero}
    \item {Almost constant}
\end{enumerate}
\vspace{2mm}

\noindent\textbf{Q11. Which logic gate produces an output of 1 only when all its inputs are 1 ?} \hfill [Topic: logic gates | Easy]
\begin{enumerate}[label=(\alph*),topsep=2pt,itemsep=1pt,partopsep=0pt,parsep=0pt]
    \item {NOR gate}
    \item {NOT gate}
    \item {OR gate}
    \item {AND gate}
\end{enumerate}
\vspace{2mm}

\noindent\textbf{Q12. A NOT gate is also commonly called a(n):} \hfill [Topic: logic gates | Easy]
\begin{enumerate}[label=(\alph*),topsep=2pt,itemsep=1pt,partopsep=0pt,parsep=0pt]
    \item {Adder}
    \item {Inverter}
    \item {Comparator}
    \item {Buffer}
\end{enumerate}
\vspace{2mm}

\noindent\textbf{Q13. The output of a two-input OR gate is 0 only when:} \hfill [Topic: logic gates | Easy]
\begin{enumerate}[label=(\alph*),topsep=2pt,itemsep=1pt,partopsep=0pt,parsep=0pt]
    \item {Both inputs are 0}
    \item {One input is 1}
    \item {Any input is 1}
    \item {Both inputs are 1}
\end{enumerate}
\vspace{2mm}

\noindent\textbf{Q14. A varactor diode is used to provide a voltage-controlled:} \hfill [Topic: Working and Operation of Varactor diode | Easy]
\begin{enumerate}[label=(\alph*),topsep=2pt,itemsep=1pt,partopsep=0pt,parsep=0pt]
    \item {Inductance}
    \item {Resistance}
    \item {Capacitance}
    \item {Frequency source}
\end{enumerate}
\vspace{2mm}

\noindent\textbf{Q15. A varactor diode is normally operated under:} \hfill [Topic: Working and Operation of Varactor diode | Easy]
\begin{enumerate}[label=(\alph*),topsep=2pt,itemsep=1pt,partopsep=0pt,parsep=0pt]
    \item {Forward bias}
    \item {Zero bias only}
    \item {No bias at all}
    \item {Reverse bias}
\end{enumerate}
\vspace{2mm}

\noindent\textbf{Q16. When testing a good diode with a multimeter, it shows low resistance in forward direction and high resistance in reverse. A diode that reads low resistance in both directions is:} \hfill [Topic: Testing of diode and BJT | Easy]
\begin{enumerate}[label=(\alph*),topsep=2pt,itemsep=1pt,partopsep=0pt,parsep=0pt]
    \item {Reverse biased}
    \item {Shorted}
    \item {Perfectly good}
    \item {Open}
\end{enumerate}
\vspace{2mm}

\noindent\textbf{Q17. A half-wave rectifier conducts during:} \hfill [Topic: Half wave and Full wave rectifier | Easy]
\begin{enumerate}[label=(\alph*),topsep=2pt,itemsep=1pt,partopsep=0pt,parsep=0pt]
    \item {Only one half of the AC cycle}
    \item {Only the negative peaks}
    \item {Both halves of the AC cycle}
    \item {The full cycle continuously}
\end{enumerate}
\vspace{2mm}

\noindent\textbf{Q18. Compared to a half-wave rectifier, a full-wave rectifier has:} \hfill [Topic: Half wave and Full wave rectifier | Easy]
\begin{enumerate}[label=(\alph*),topsep=2pt,itemsep=1pt,partopsep=0pt,parsep=0pt]
    \item {Higher output efficiency}
    \item {The same single diode}
    \item {A}
    \item {No output voltage}
\end{enumerate}
\vspace{2mm}

\noindent\textbf{Q19. How many diodes are used in a full-wave bridge rectifier?} \hfill [Topic: Full wave bridge rectifier | Easy]
\begin{enumerate}[label=(\alph*),topsep=2pt,itemsep=1pt,partopsep=0pt,parsep=0pt]
    \item {Two}
    \item {Four}
    \item {One}
    \item {Six}
\end{enumerate}
\vspace{2mm}

\noindent\textbf{Q20. What does the abbreviation MOSFET stand for?} \hfill [Topic: MOSFET representation and characteristics | Easy]
\begin{enumerate}[label=(\alph*),topsep=2pt,itemsep=1pt,partopsep=0pt,parsep=0pt]
    \item {Metal Only Semiconductor Fast Emitting Transistor}
    \item {Multiple Output Switching Field Effect Transistor}
    \item {Metal Oxide Semiconductor Field Effect Transistor}
    \item {Metal Oxide Silicon Forward Emitter Transistor}
\end{enumerate}
\vspace{2mm}

\noindent\textbf{Q21. A pure silicon crystal is doped with a pentavalent impurity such as phosphorus. What type of semiconductor is formed and what are the majority carriers?} \hfill [Topic: introduction to semiconductor devices | Medium]
\begin{enumerate}[label=(\alph*),topsep=2pt,itemsep=1pt,partopsep=0pt,parsep=0pt]
    \item {Intrinsic semiconductor with equal electrons and holes}
    \item {N-type with electrons as majority carriers}
    \item {P-type with holes as majority carriers}
    \item {N-type with holes as majority carriers}
\end{enumerate}
\vspace{2mm}

\noindent\textbf{Q22. When a PN junction diode is forward biased, what happens to the width of the depletion region and the potential barrier?} \hfill [Topic: Working and operation of PN Junction diode | Medium]
\begin{enumerate}[label=(\alph*),topsep=2pt,itemsep=1pt,partopsep=0pt,parsep=0pt]
    \item {Both increase, blocking current flow}
    \item {Width increases but barrier decreases}
    \item {Both decrease, allowing current to flow}
    \item {Width decreases but barrier increases}
\end{enumerate}
\vspace{2mm}

\noindent\textbf{Q23. In the V-I characteristics of a silicon PN junction diode, appreciable forward current begins only after the applied voltage exceeds approximately which value?} \hfill [Topic: Characteristics of PN junction diode | Medium]
\begin{enumerate}[label=(\alph*),topsep=2pt,itemsep=1pt,partopsep=0pt,parsep=0pt]
    \item {1.1 V}
    \item {0.0 V}
    \item {0.7 V}
    \item {0.3 V}
\end{enumerate}
\vspace{2mm}

\noindent\textbf{Q24. A Zener diode used for voltage regulation is always connected in which configuration in the circuit?} \hfill [Topic: working and operation of Zener diode | Medium]
\begin{enumerate}[label=(\alph*),topsep=2pt,itemsep=1pt,partopsep=0pt,parsep=0pt]
    \item {Reverse biased and operated in the breakdown region}
    \item {Reverse biased but operated below breakdown}
    \item {Forward biased near its knee voltage}
    \item {Forward biased and operated in saturation}
\end{enumerate}
\vspace{2mm}

\noindent\textbf{Q25. In the reverse breakdown region of a Zener diode's characteristic curve, how do voltage and current behave?} \hfill [Topic: characteristics of Zener diode | Medium]
\begin{enumerate}[label=(\alph*),topsep=2pt,itemsep=1pt,partopsep=0pt,parsep=0pt]
    \item {Voltage increases rapidly while current stays constant}
    \item {Both voltage and current decrease together}
    \item {Voltage remains almost constant while current changes significantly}
    \item {Both voltage and current increase proportionally}
\end{enumerate}
\vspace{2mm}

\noindent\textbf{Q26. A two-input logic gate produces a HIGH output only when both inputs are LOW. Which gate is this?} \hfill [Topic: logic gates | Medium]
\begin{enumerate}[label=(\alph*),topsep=2pt,itemsep=1pt,partopsep=0pt,parsep=0pt]
    \item {XOR gate}
    \item {AND gate}
    \item {NOR gate}
    \item {NAND gate}
\end{enumerate}
\vspace{2mm}

\noindent\textbf{Q27. Using De Morgan's theorem, the expression A $\cdot$ B is equivalent to which of the following?} \hfill [Topic: logic gates | Medium]
\begin{enumerate}[label=(\alph*),topsep=2pt,itemsep=1pt,partopsep=0pt,parsep=0pt]
    \item {A + B}
    \item {A $\cdot$ B}
    \item {A $\cdot$ B}
    \item {A + B}
\end{enumerate}
\vspace{2mm}

\noindent\textbf{Q28. A varactor diode is operated under reverse bias so that it behaves like a voltage-controlled capacitor. As the reverse voltage increases, what happens to its capacitance?} \hfill [Topic: Working and Operation of Varactor diode | Medium]
\begin{enumerate}[label=(\alph*),topsep=2pt,itemsep=1pt,partopsep=0pt,parsep=0pt]
    \item {Capacitance increases because the depletion width increases}
    \item {Capacitance decreases because the depletion width increases}
    \item {Capacitance stays constant regardless of reverse voltage}
    \item {Capacitance increases because the depletion width decreases}
\end{enumerate}
\vspace{2mm}

\noindent\textbf{Q29. While testing a diode with a digital multimeter in diode mode, a technician reads about 0.7 V in one direction and "OL" (open) when the leads are reversed. What is the condition of the diode?} \hfill [Topic: Testing of diode and BJT | Medium]
\begin{enumerate}[label=(\alph*),topsep=2pt,itemsep=1pt,partopsep=0pt,parsep=0pt]
    \item {The diode is leaky}
    \item {The diode is shorted}
    \item {The diode is open in both directions}
    \item {The diode is good}
\end{enumerate}
\vspace{2mm}

\noindent\textbf{Q30. When testing an NPN transistor with a multimeter, the base-emitter and base-collector junctions should behave like which arrangement of diodes?} \hfill [Topic: Testing of diode and BJT | Medium]
\begin{enumerate}[label=(\alph*),topsep=2pt,itemsep=1pt,partopsep=0pt,parsep=0pt]
    \item {Two diodes with their anodes joined at the base}
    \item {Two diodes with their cathodes joined at the base}
    \item {Two diodes connected in series between collector and emitter}
    \item {A single diode between collector and emitter}
\end{enumerate}
\vspace{2mm}

\noindent\textbf{Q31. A half-wave rectifier is supplied with a 50 Hz AC input. What is the ripple frequency of its output?} \hfill [Topic: Half wave and Full wave rectifier | Medium]
\begin{enumerate}[label=(\alph*),topsep=2pt,itemsep=1pt,partopsep=0pt,parsep=0pt]
    \item {100 Hz}
    \item {150 Hz}
    \item {50 Hz}
    \item {25 Hz}
\end{enumerate}
\vspace{2mm}

\noindent\textbf{Q32. A full-wave rectifier produces a peak output voltage of 10 V . What is the approximate average (DC) output voltage?} \hfill [Topic: Half wave and Full wave rectifier | Medium]
\begin{enumerate}[label=(\alph*),topsep=2pt,itemsep=1pt,partopsep=0pt,parsep=0pt]
    \item {10 V}
    \item {6.37 V}
    \item {7.07 V}
    \item {3.18 V}
\end{enumerate}
\vspace{2mm}

\noindent\textbf{Q33. What is the maximum theoretical rectification efficiency of a full-wave rectifier?} \hfill [Topic: Half wave and Full wave rectifier | Medium]
\begin{enumerate}[label=(\alph*),topsep=2pt,itemsep=1pt,partopsep=0pt,parsep=0pt]
    \item {50\%}
    \item {40.6\%}
    \item {81.2\%}
    \item {100\%}
\end{enumerate}
\vspace{2mm}

\noindent\textbf{Q34. In a full-wave bridge rectifier using a single transformer secondary of peak voltage V m , approximately what is the Peak Inverse Voltage (PIV) across each diode?} \hfill [Topic: Full wave bridge rectifier | Medium]
\begin{enumerate}[label=(\alph*),topsep=2pt,itemsep=1pt,partopsep=0pt,parsep=0pt]
    \item {0.5 V m}
    \item {2 V m}
    \item {V m}
    \item {V m /2}
\end{enumerate}
\vspace{2mm}

\noindent\textbf{Q35. A key advantage of a bridge rectifier over a center-tapped full-wave rectifier is that it:} \hfill [Topic: Full wave bridge rectifier | Medium]
\begin{enumerate}[label=(\alph*),topsep=2pt,itemsep=1pt,partopsep=0pt,parsep=0pt]
    \item {Produces a higher ripple frequency}
    \item {Requires no transformer at all}
    \item {Does not require a center-tapped transformer}
    \item {Uses only one diode}
\end{enumerate}
\vspace{2mm}

\noindent\textbf{Q36. An enhancement-type N-channel MOSFET starts conducting between drain and source only when the gate-to-source voltage satisfies which condition?} \hfill [Topic: MOSFET representation and characteristics | Medium]
\begin{enumerate}[label=(\alph*),topsep=2pt,itemsep=1pt,partopsep=0pt,parsep=0pt]
    \item {V GS < 0}
    \item {V GS = V D S}
    \item {V GS > V t h (threshold voltage)}
    \item {V GS = 0}
\end{enumerate}
\vspace{2mm}

\noindent\textbf{Q37. In a MOSFET, the gate terminal controls the drain current but draws almost no DC current because:} \hfill [Topic: MOSFET representation and characteristics | Medium]
\begin{enumerate}[label=(\alph*),topsep=2pt,itemsep=1pt,partopsep=0pt,parsep=0pt]
    \item {The gate forms a forward-biased PN junction}
    \item {The gate is insulated from the channel by a thin oxide layer}
    \item {The gate is made of a highly resistive metal}
    \item {The gate is directly connected to the source}
\end{enumerate}
\vspace{2mm}

\noindent\textbf{Q38. In a power supply circuit, a capacitor is connected across the output of a rectifier. Its primary function is to:} \hfill [Topic: Application of diodes, transistors, rectifiers and MOSFET | Medium]
\begin{enumerate}[label=(\alph*),topsep=2pt,itemsep=1pt,partopsep=0pt,parsep=0pt]
    \item {Convert AC into DC}
    \item {Rectify the negative half-cycles}
    \item {Smooth the output by reducing ripple}
    \item {Increase the peak output voltage indefinitely}
\end{enumerate}
\vspace{2mm}

\noindent\textbf{Q39. A diode is placed across the coil of a relay in a circuit driven by a transistor. The purpose of this diode is to:} \hfill [Topic: Application of diodes, transistors, rectifiers and MOSFET | Medium]
\begin{enumerate}[label=(\alph*),topsep=2pt,itemsep=1pt,partopsep=0pt,parsep=0pt]
    \item {Act as a voltage regulator for the coil}
    \item {Protect the transistor by suppressing the inductive voltage spike when the coil switches off}
    \item {Increase the current flowing through the relay coil}
    \item {Rectify the AC supply feeding the relay}
\end{enumerate}
\vspace{2mm}

\noindent\textbf{Q40. The reverse saturation current of a PN junction diode is primarily due to:} \hfill [Topic: Characteristics of PN junction diode | Medium]
\begin{enumerate}[label=(\alph*),topsep=2pt,itemsep=1pt,partopsep=0pt,parsep=0pt]
    \item {The applied forward voltage}
    \item {Minority carriers crossing the junction}
    \item {The barrier potential increasing}
    \item {Majority carriers crossing the junction}
\end{enumerate}
\vspace{2mm}

\noindent\textbf{Q41. A silicon diode has a reverse saturation current I S = 10 nA at 300 K . Using the ideal diode equation with V T = 25.9 mV and ideality factor n = 1 , what forward voltage is required to produce a current of 1 mA ?} \hfill [Topic: Characteristics of PN junction diode | Hard]
\begin{enumerate}[label=(\alph*),topsep=2pt,itemsep=1pt,partopsep=0pt,parsep=0pt]
    \item {Approximately 0.477 V}
    \item {Approximately 0.600 V}
    \item {Approximately 0.298 V}
    \item {Approximately 0.700 V}
\end{enumerate}
\vspace{2mm}

\noindent\textbf{Q42. In a PN junction at equilibrium, the built-in potential barrier opposes diffusion of majority carriers. If the doping on both sides is increased by a factor of 10, how does the built-in potential V bi change (temperature constant)?} \hfill [Topic: Working and operation of PN Junction diode | Hard]
\begin{enumerate}[label=(\alph*),topsep=2pt,itemsep=1pt,partopsep=0pt,parsep=0pt]
    \item {It increases by V T ln 10}
    \item {It increases by 2 V T ln 10}
    \item {It remains unchanged}
    \item {It decreases by 2 V T ln 10}
\end{enumerate}
\vspace{2mm}

\noindent\textbf{Q43. A Zener regulator uses a 6.2 V Zener with dynamic resistance r Z = 5 $\Omega$ . If the Zener current changes from 5 mA to 25 mA , what is the change in output (Zener) voltage?} \hfill [Topic: characteristics of Zener diode | Hard]
\begin{enumerate}[label=(\alph*),topsep=2pt,itemsep=1pt,partopsep=0pt,parsep=0pt]
    \item {0.10 V}
    \item {0.05 V}
    \item {0.125 V}
    \item {0.25 V}
\end{enumerate}
\vspace{2mm}

\noindent\textbf{Q44. In a shunt Zener regulator with V in = 15 V , series resistor R S = 220 $\Omega$ , and a 9 V Zener, what is the maximum load current the circuit can supply while keeping the Zener in regulation (ignore minimum Zener current)?} \hfill [Topic: working and operation of Zener diode | Hard]
\begin{enumerate}[label=(\alph*),topsep=2pt,itemsep=1pt,partopsep=0pt,parsep=0pt]
    \item {Approximately 27.3 mA}
    \item {Approximately 68.2 mA}
    \item {Approximately 40.9 mA}
    \item {Approximately 18.0 mA}
\end{enumerate}
\vspace{2mm}

\noindent\textbf{Q45. A half-wave rectifier delivers a peak output of 20 V into a resistive load. What is the ripple factor of the output?} \hfill [Topic: Half wave and Full wave rectifier | Hard]
\begin{enumerate}[label=(\alph*),topsep=2pt,itemsep=1pt,partopsep=0pt,parsep=0pt]
    \item {0.318}
    \item {0.636}
    \item {0.482}
    \item {1.21}
\end{enumerate}
\vspace{2mm}

\noindent\textbf{Q46. A full-wave bridge rectifier is fed from a transformer secondary with RMS voltage 18 V . Assuming silicon diodes with 0.7 V drop each, what is the approximate DC output voltage (no filter)?} \hfill [Topic: Full wave bridge rectifier | Hard]
\begin{enumerate}[label=(\alph*),topsep=2pt,itemsep=1pt,partopsep=0pt,parsep=0pt]
    \item {Approximately 16.2 V}
    \item {Approximately 7.6 V}
    \item {Approximately 11.5 V}
    \item {Approximately 15.3 V}
\end{enumerate}
\vspace{2mm}

\noindent\textbf{Q47. For a full-wave bridge rectifier, how does the Peak Inverse Voltage (PIV) rating per diode compare to a center-tapped full-wave rectifier producing the same peak output V m ?} \hfill [Topic: Full wave bridge rectifier | Hard]
\begin{enumerate}[label=(\alph*),topsep=2pt,itemsep=1pt,partopsep=0pt,parsep=0pt]
    \item {Bridge PIV = V m , center-tapped PIV = 2 V m}
    \item {Both have PIV = V m}
    \item {Both have PIV = 2 V m}
    \item {Bridge PIV = 2 V m , center-tapped PIV = V m}
\end{enumerate}
\vspace{2mm}

\noindent\textbf{Q48. A varactor diode has junction capacitance C j = C 0 / ( 1 + V R / V bi ) 1/2 with C 0 = 80 pF and V bi = 0.7 V . What is C j at a reverse bias of V R = 8.3 V ?} \hfill [Topic: Working and Operation of Varactor diode | Hard]
\begin{enumerate}[label=(\alph*),topsep=2pt,itemsep=1pt,partopsep=0pt,parsep=0pt]
    \item {40.0 pF}
    \item {24.6 pF}
    \item {56.6 pF}
    \item {8.0 pF}
\end{enumerate}
\vspace{2mm}

\noindent\textbf{Q49. A logic function is given by F = A + B + A $\cdot$ B . Which single gate is equivalent to F ?} \hfill [Topic: logic gates | Hard]
\begin{enumerate}[label=(\alph*),topsep=2pt,itemsep=1pt,partopsep=0pt,parsep=0pt]
    \item {A NOR gate}
    \item {An XOR gate}
    \item {An AND gate}
    \item {A NOT gate on A (i.e., F = A )}
\end{enumerate}
\vspace{2mm}

\noindent\textbf{Q50. Using only 2-input NAND gates, what is the minimum number of NAND gates required to implement a 2-input XOR function?} \hfill [Topic: logic gates | Hard]
\begin{enumerate}[label=(\alph*),topsep=2pt,itemsep=1pt,partopsep=0pt,parsep=0pt]
    \item {3}
    \item {5}
    \item {4}
    \item {2}
\end{enumerate}
\vspace{2mm}

\noindent\textbf{Q51. An n-channel enhancement MOSFET has V T H = 1.5 V and k = 0.5 mA/V 2 (where I D = k ( V GS - V T H ) 2 in saturation). For V GS = 3.5 V and V D S = 5 V , what is I D ?} \hfill [Topic: MOSFET representation and characteristics | Hard]
\begin{enumerate}[label=(\alph*),topsep=2pt,itemsep=1pt,partopsep=0pt,parsep=0pt]
    \item {4.0 mA}
    \item {1.0 mA}
    \item {2.0 mA}
    \item {8.0 mA}
\end{enumerate}
\vspace{2mm}

\noindent\textbf{Q52. An n-channel MOSFET with V T H = 1 V operates with V GS = 4 V . At what maximum V D S does the device remain in the triode (ohmic) region?} \hfill [Topic: MOSFET representation and characteristics | Hard]
\begin{enumerate}[label=(\alph*),topsep=2pt,itemsep=1pt,partopsep=0pt,parsep=0pt]
    \item {Below 5 V}
    \item {Below 1 V}
    \item {Below 3 V}
    \item {Below 4 V}
\end{enumerate}
\vspace{2mm}

\noindent\textbf{Q53. Using a digital multimeter in diode-test mode on a good silicon NPN transistor, which set of readings is expected?} \hfill [Topic: Testing of diode and BJT | Hard]
\begin{enumerate}[label=(\alph*),topsep=2pt,itemsep=1pt,partopsep=0pt,parsep=0pt]
    \item {All three junctions read $\approx$ 0 V in both directions}
    \item {Forward B-E reads 0.3 V and B-C reads open}
    \item {B-E and B-C read open in both directions; C-E reads 0.7 V}
    \item {Forward B-E and B-C both read $\approx$ 0.6 -- 0.7 V ; C-E reads open both ways}
\end{enumerate}
\vspace{2mm}

\noindent\textbf{Q54. A diode under test shows a low resistance reading (near 0 $\Omega$ ) in BOTH forward and reverse directions with an ohmmeter. What is the most likely condition?} \hfill [Topic: Testing of diode and BJT | Hard]
\begin{enumerate}[label=(\alph*),topsep=2pt,itemsep=1pt,partopsep=0pt,parsep=0pt]
    \item {The diode is shorted}
    \item {The diode is leaky but usable}
    \item {The diode is healthy}
    \item {The diode is open}
\end{enumerate}
\vspace{2mm}

\noindent\textbf{Q55. The reverse saturation current of a silicon diode approximately doubles for every $10^\circ\text{C}$ rise. If I S = 5 nA at $25^\circ\text{C}$ , what is I S at $65^\circ\text{C}$ ?} \hfill [Topic: Characteristics of PN junction diode | Hard]
\begin{enumerate}[label=(\alph*),topsep=2pt,itemsep=1pt,partopsep=0pt,parsep=0pt]
    \item {40 nA}
    \item {80 nA}
    \item {20 nA}
    \item {160 nA}
\end{enumerate}
\vspace{2mm}

\noindent\textbf{Q56. A full-wave rectifier (center-tapped) and a half-wave rectifier both produce the same peak voltage V m . What is the ratio of their DC output voltages (full-wave to half-wave)?} \hfill [Topic: Half wave and Full wave rectifier | Hard]
\begin{enumerate}[label=(\alph*),topsep=2pt,itemsep=1pt,partopsep=0pt,parsep=0pt]
    \item {2 : 1}
    \item {1 : 1}
    \item {$\pi$ : 1}
    \item {2 : 1}
\end{enumerate}
\vspace{2mm}

\noindent\textbf{Q57. In a diode clamper circuit, a capacitor charges to the peak input voltage and shifts the DC level of the signal. For a positive clamper with input $\pm$ 10 V sine and an ideal diode, what are the output voltage limits?} \hfill [Topic: Application of diodes, transistors, rectifiers and MOSFET | Hard]
\begin{enumerate}[label=(\alph*),topsep=2pt,itemsep=1pt,partopsep=0pt,parsep=0pt]
    \item {- 10 V to + 10 V}
    \item {- 20 V to 0 V}
    \item {0 V to + 10 V}
    \item {0 V to + 20 V}
\end{enumerate}
\vspace{2mm}

\noindent\textbf{Q58. Intrinsic silicon has n i = 1.5 $\times$ 1 0 10 cm - 3 . If it is doped with donors at N D = 1.5 $\times$ 1 0 16 cm - 3 , what is the approximate equilibrium hole concentration p ?} \hfill [Topic: introduction to semiconductor devices | Hard]
\begin{enumerate}[label=(\alph*),topsep=2pt,itemsep=1pt,partopsep=0pt,parsep=0pt]
    \item {1.5 $\times$ 1 0 16 cm - 3}
    \item {1.0 $\times$ 1 0 6 cm - 3}
    \item {1.5 $\times$ 1 0 10 cm - 3}
    \item {1.5 $\times$ 1 0 4 cm - 3}
\end{enumerate}
\vspace{2mm}

\noindent\textbf{Q59. Below approximately 5 V , Zener diodes break down mainly by one mechanism, and above approximately 6 V by another. Which statement correctly describes their temperature coefficients?} \hfill [Topic: working and operation of Zener diode | Hard]
\begin{enumerate}[label=(\alph*),topsep=2pt,itemsep=1pt,partopsep=0pt,parsep=0pt]
    \item {Both mechanisms have positive temperature coefficients}
    \item {Both mechanisms have negative temperature coefficients}
    \item {Zener breakdown ( < 5 V ) has negative tempco; avalanche ( > 6 V ) has positive tempco}
    \item {Zener breakdown has positive tempco; avalanche has negative tempco}
\end{enumerate}
\vspace{2mm}

\noindent\textbf{Q60. A capacitor-filtered full-wave rectifier supplies a load drawing I L = 100 mA from a 1000 $\mu$ F capacitor at a line frequency of 50 Hz . What is the approximate peak-to-peak ripple voltage?} \hfill [Topic: Application of diodes, transistors, rectifiers and MOSFET | Hard]
\begin{enumerate}[label=(\alph*),topsep=2pt,itemsep=1pt,partopsep=0pt,parsep=0pt]
    \item {2.0 V}
    \item {0.5 V}
    \item {1.0 V}
    \item {5.0 V}
\end{enumerate}
\vspace{2mm}


\newpage
\section{Section B: Comprehensive Analytical Problems with Step-by-Step Solutions}

\begin{subjectiveq}{4}{1}
A bridge rectifier is supplied from a $230\text{ V}$, $50\text{ Hz}$ AC mains via a step-down transformer having a turns ratio of $10:1$. The diodes have forward dynamic resistance $r_d = 0.5\,\Omega$ each and the load resistor is $R_L = 100\,\Omega$. Calculate:\begin{enumerate}[label=(\alph*)]\item The peak secondary voltage $V_m$, DC load current $I_{dc}$, and DC output voltage $V_{dc}$.\item Rectification efficiency $\eta$ and ripple factor $\gamma$ without a filter.\item If a shunt capacitor filter of $C = 1000\,\mu\text{F}$ is connected across the load, determine the resulting ripple factor and DC output voltage.\end{enumerate}
\end{subjectiveq}

\begin{solbox}
\textbf{Step 1: Secondary Voltage and DC Output Quantities}\\RMS secondary voltage: $V_{2,rms} = \frac{230}{10} = 23\text{ V}$.Peak secondary voltage: $V_m = 23\sqrt{2} \approx 32.53\text{ V}$.Two conducting diodes in series per half-cycle ($2r_d = 1.0\,\Omega$):\[ I_m = \frac{V_m}{2r_d + R_L} = \frac{32.53}{1.0 + 100} = \frac{32.53}{101} \approx 0.322\text{ A} \]DC load current: $I_{dc} = \frac{2 I_m}{\pi} \approx 0.205\text{ A}$.DC output voltage: $V_{dc} = I_{dc} R_L = 0.205 \times 100 = 20.5\text{ V}$.\textbf{Step 2: Rectification Efficiency and Unfiltered Ripple Factor}\\Rectification efficiency: $\eta = \frac{8}{\pi^2}\frac{R_L}{2r_d + R_L} \approx 80.4\%$.Ripple factor without filter: $\gamma = \sqrt{\left(\frac{\pi}{2\sqrt{2}}\right)^2 - 1} \approx 0.482$ ($48.2\%$).\textbf{Step 3: Capacitor Filter Performance}\\Ripple factor with shunt capacitor filter ($f = 50\text{ Hz}$):\[ \gamma = \frac{1}{4\sqrt{3} f C R_L} = \frac{1}{4\sqrt{3} \times 50 \times (1000 \times 10^{-6}) \times 100} \approx 0.0289 \text{ (} 2.89\% \text{)} \]DC voltage under capacitive filtering: $V_{dc} = V_m - \frac{I_{dc}}{4 f C} \approx 31.51\text{ V}$.
\end{solbox}

\begin{examinertip}
\textbf{Examiner Insight}: In bridge rectifiers, PIV per diode is $V_m$, which is half the $2V_m$ required in center-tapped transformers.
\end{examinertip}
\vspace{4mm}
